\medskip

\emph{Principaux domaines abordés}: Étude des fonctions.
Fonction logarithme.

\bigskip

\textbf{Partie A}

\begin{center}
\psset{unit=0.6cm,arrowsize=2pt 3}
\begin{pspicture*}(-4,-5)(15.5,8)
\psgrid[gridlabels=0pt,subgriddiv=1,gridwidth=0.15pt]
\psaxes[linewidth=1.25pt,labelFontSize=\scriptstyle]{->}(0,0)(-4,-5)(15.5,8)
\psplot[plotpoints=2000,linewidth=1.25pt,linecolor=blue]{3.001}{15.5}{x dup mul x sub 6 sub ln}
\psplot[plotpoints=2000,linewidth=1.25pt,linecolor=red,linestyle=dashed]{3.001}{15.5}{x 2 mul 1 sub x dup mul x sub 6 sub div}
\uput[r](3.1,7.7){\red $\mathcal{C}_2$}\uput[r](3.1,-4){\blue $\mathcal{C}_1$}
\psline[ArrowInside=->](0,3)(5.63,3)(5.62,0)
\end{pspicture*}
\end{center}

Dans le repère orthonormé ci-dessus, sont tracées les courbes représentatives d'une fonction $f$ et de sa fonction dérivée, notée $f'$, toutes deux définies sur $]3~;~+\infty[$.

\medskip

\begin{enumerate}
\item 
Si $\mathcal{C}_1$ était la représentation de la dérivée, celle-ci serait négative puis positive , donc la fonction serait décroissante puis croissante ce qui n'est pas le cas de la fonction représente par la courbe  $\mathcal{C}_2$.

Donc $\mathcal{C}_2$ est la représentation de la fonction dérivée et $\mathcal{C}_1$ celle de la fonction.
\item 
On lit approximativement $f(5,6) \approx 3$. 
\item 
$f$ semble concave sur $]3~;~+\infty[$.
\end{enumerate}

\bigskip

\textbf{Partie B}

\medskip

\begin{enumerate}
\item 
$\bullet~~$ $x^2 - x - 6 = \left(x - \frac12\right)^2 - \frac14 - 6 = \left(x - \frac12\right)^2 - \frac{25}{4} = \left(x - \frac12\right)^2 - \left(\frac{5}{2}\right)^2 = \left(x - \frac12 + \frac52 \right)\left(x - \frac12 - \frac52 \right) = (x + 2)(x - 3)$.

On sait que ce trinôme est positif sauf sur l'intervalle $[-2~;~3]$. Donc $\ln \left(x^2- x- 6\right)$ est bien définie sur $I = ]3~;~+ \infty[$ ;

$\bullet~~$On peut calculer $\Delta = 1 + 24 = 25 = 5^2$ et en déduire les racines du trinôme $x^2- x- 6$ : $x_1 = \dfrac{1 + 5}{2} = 3$ et $x_2 = \dfrac{1-5}{2} = - 2$ et faire ensuite la même conclusion.
\item On admet que la fonction $f$ de la Partie A est définie par $f(x) = \ln \left(x^2- x- 6\right)$ sur $I$.




$\bullet~~$Limite en $+ \infty$  : $\ln \left(x^2- x- 6\right) = \ln x^2\left(1 - \dfrac{1}{x} - \dfrac{6}{x^2} \right)$.

On a $\displaystyle\lim_{x \to + \infty} \dfrac{1}{x} = \displaystyle\lim_{x \to + \infty} \dfrac{6}{x^2} = 0$, donc 

$\displaystyle\lim_{x \to + \infty} 1 - \dfrac{1}{x} - \dfrac{6}{x^2} = 1$.

Comme $\displaystyle\lim_{x \to + \infty} x^2 = + \infty$, par composition de la fonction logarithme $\displaystyle\lim_{x \to + \infty} f(x) = + \infty$.

$\bullet~~$Limite en 3 : $f(x) = \ln (x + 2)(x - 3)$.

On a $\displaystyle\lim_{x \to + 3}(x + 2)(x - 3) = 5 \times 0 = 0$, donc $\displaystyle\lim_{x \to + 3} \ln (x + 2)(x - 3) = - \infty$. Ceci signifie que la droite verticale d'équation $x = 3$ est asymptote à courbe représentative de la fonction $f$ sur $I$ (soit $\mathcal{C}_1$)
\item 
	\begin{enumerate}
		\item 
		Sur $I$, on a vu que $x^2 - x - 6 > 0$, donc la fonction $f$ est dérivable et sur $I$, on a $f'(x) = \dfrac{2x - 1}{x^2 - x - 6}$.
		\item 
$f'(x)$ a le signe du numérateur $2x - 1$ :

pour $x>3$, on a $2x-1>5$ donc $2x-1>0$ et donc $f'(x>0$.

Donc la fonction $f$ est strictement croissante sur $I$ de moins l'infini à plus l'infini.	

On dresse le tableau des variations de la fonction $f$.

\begin{center}
{\renewcommand{\arraystretch}{1.3}
     
\psset{nodesep=3pt,arrowsize=2pt 3}
     
\def\esp{\hspace*{3cm}}
     
\def\hauteur{0pt}
     
$\begin{array}{|c| l *2{c}|}
\hline
 x & 3   & \esp & +\infty \\
 \hline
f'(x) & \vline\;\vline  & \pmb{+} & \\  
\hline
  & \vline\;\vline  &    & \Rnode{max}{+\infty}   \\
f(x) & \vline\;\vline&  &  \rule{0pt}{\hauteur} \\
 &\vline\;\vline     \Rnode{min}{~-\infty} & & \rule{0pt}{\hauteur}
\ncline{->}{min}{max}\\
\hline
\end{array}$
}
\end{center}


	\end{enumerate}
\item 
	\begin{enumerate}
		\item 
		La fonction $f$ est strictement croissante sur $I$ et en particulier sur ]5~;~6[.
		
De plus $f(5) = \ln (25 - 5 - 6) = \ln 14 \approx 2,64<3$ et $f(6) = \ln (36 - 6 - 6) = \ln 24 \approx 3,18>3$.

La fonction $f$ étant continue car dérivable sur l'intervalle ]5~;~6[, d'après la propriété des valeurs intermédiaires il existe un réel unique $\alpha$ tel que $f(\alpha) = 3$, avec $5 < \alpha < 6$.

		\item 
		La calculatrice donne :
		
$f(5,6) \approx 2,984$ et $f(5,7) \approx 3,035$, donc $5,6 < \alpha < 5,7$ puis 

$f(5,63) \approx 2,999$ et $f(5,64) \approx 3,004$ donc $5,63 < \alpha < 5,64$.
	\end{enumerate}
\item 
	\begin{enumerate}
		\item 
		La fonction $f'$ quotient de fonctions dérivables sur $I$, le dénominateur ne s'annulant pas est elle-même dérivable sur $I$ et sur cet intervalle :
		
		$f''(x) = \dfrac{2\left (x^2 - x - 6\right) - (2x - 1)(2x - 1)}{\left (x^2 - x - 6\right)^2} = \dfrac{2x^2  - 2x - 12 - 4x^2 - 1 - 4x}{\left (x^2 - x - 6\right)^2}  = \dfrac{- 2x^2 + 2x - 13}{\left(x^2 - x - 6\right)^2}$.
		\item 
Le signe de la dérivée seconde est celui du numérateur puisque $\left(x^2 - x - 6\right)^2 > 0$.

Or ce trinôme $- 2x^2 + 2x - 13$ n'a pas de racine ($\Delta = 4 - 8 \times 13 = - 100 < 0$), donc est négatif en particulier sur $I$.

Sur $I$, $f''(x) < 0$ : la fonction est concave sur cet intervalle.
	\end{enumerate}
\end{enumerate}
